% Manuscript subsection for the ECON-08 systems paper. Tables use only standard LaTeX.
\subsection{Protocol deviations}

The access protocol specified one Dune query returning only aggregate interval bins and counts. The query returned HTTP 402. We subsequently queried public Aave v3 Base Supply and Withdraw logs by RPC and reconstructed access sequences locally. This exploratory change preserved the order of events within an address. It allowed a paired comparison of cache policies along the same observed path. The published protocol, the unsuccessful Dune request and this deviation are retained in the research record.

We requested 180 days and obtained 58 consecutive UTC days, 8 September through 4 November 2026. We selected the longest fully covered span using data availability before computing interval statistics. We used the beneficiary for supplies and the owner for withdrawals. Multiple events concerning the same address in one transaction counted once. The window contained 42,739 distinct addresses. Of these, 14,440 made at least three accesses and entered the analysis. Their 362,656 accesses gave 348,216 adjacent intervals. Intervals are censored at the boundaries. All cost figures below refer to this active cohort over 58 days.

The local input contains public addresses, transaction hashes and timestamps. No association between an Aave address and an EAS issuer was observed or inferred. Individual access rows and historical account-code lookup tables remain outside the public repository. From 18 November the retained local copies used owner-only directories and files. A chain node received the public addresses for bytecode queries. The public release contains aggregate tables and block-level provenance. The original protocol set no fixed retention period for the fallback inputs.

\subsection{Access concentration and account-code sensitivity}

Before computing the subgroup results we recorded an analysis addendum. We queried account bytecode at closing Base block 53,523,726. Among the eligible addresses, 9,998 returned empty bytecode, 3,060 had an EIP-7702 delegation designation, 1,382 had other nonempty bytecode and 0 had an unknown lookup outcome. Empty code at the end of the window is a proxy for a conventional externally owned account (EOA). Code may change over time and an EOA can execute automated strategies. We applied daily caps to complete addresses. An address was excluded if its number of distinct transactions exceeded the cap on any UTC day. Activity quartiles used total transactions over the window and kept identical counts together.

\begin{center}
\small
\begin{tabular}{lrrrr}
\hline
Cohort & Addresses & Pooled gaps & Equal wallet & Extra proofs \\
 & & below 1 h (\%) & below 1 h (\%) & at $F=1$ h \\
\hline
All eligible & 14,440 & 77.13 & 38.40 & 2.66 \\
Empty code & 9,998 & 51.88 & 40.46 & 1.33 \\
Delegated EOA & 3,060 & 40.27 & 33.19 & 1.77 \\
Other code & 1,382 & 90.21 & 35.11 & 14.22 \\
Cap of 10 per day & 14,047 & 35.15 & 37.54 & 1.17 \\
Empty code and cap of 10 & 9,832 & 38.20 & 39.97 & 0.93 \\
Cap of 25 per day & 14,322 & 38.05 & 38.06 & 1.50 \\
Cap of 50 per day & 14,372 & 39.89 & 38.17 & 1.75 \\
\hline
\end{tabular}
\end{center}

Pooled percentages weight each interval equally. Equal-wallet percentages first compute the proportion for each address and then average over addresses. The extra-proof column gives the common-epoch mean minus the selective mean for two equally weighted EAS sources. At the cap of ten accesses per day, 393 addresses were excluded. Those addresses accounted for 72.06\% of all eligible accesses. The 1,382 other code-bearing accounts generated 65.91\% of all eligible accesses. The difference between 77.13\% of pooled intervals and 38.40\% under equal wallet weights is therefore consequential. The cost increase persists for conventional empty-code accounts after this cap, at 0.93 proofs per address. Activity strata contain ties. Their observed ranges are three, four, five to eight and at least nine transactions, respectively. At a one-hour freshness ceiling, the four respective increments were 0.20, 0.31, 0.73 and 10.76 proofs per address.

\subsection{Cost and dependence of the cache decision on the issuer}

Each observed access path was evaluated under each of the two hypothetical EAS source schedules with equal source weights, a cold cache, a 30-day time to live and five minutes of assumed publication and confirmation lag. For each access we recorded whether switching the source schedule changed the full-proof versus cache-hit decision. The same-time comparison holds accesses, deadlines and cache rules fixed. It is a measure of dependence on the issuer schedule. It gives no estimate of longitudinal issuer-identification accuracy.

\begin{center}
\small
\begin{tabular}{rrrrrr}
\hline
$F$ (h) & Selective & Shared & Fixed batch & Different decisions & Different decisions \\
 & \multicolumn{3}{c}{full proofs per wallet} & per access (\%) & per wallet (\%) \\
\hline
1 & 5.72 & 8.38 & 8.39 & 21.18 & 21.95 \\
6 & 4.49 & 5.93 & 5.93 & 11.45 & 17.33 \\
24 & 3.91 & 4.76 & 4.76 & 6.80 & 13.03 \\
\hline
\end{tabular}
\end{center}

The last two columns apply to selective publication. In every cohort and at all three deadlines, each change was from a cache hit under the intermittent source to a full proof under the daily source. The pooled figure is thus the difference in the two full-proof totals divided by accesses. It provides no independent evidence of successful identification. For a shared epoch or a fixed batch both sources produce the same decision on every matched access. With equal source weights, chance accuracy for the common-epoch assignment is one half by symmetry. Only group sizes one and two are supported by these two measured schedules. The larger grouping study uses replicated templates and is a separate model experiment. At the one-hour freshness ceiling, holders assigned to the intermittent source incur 5.32 additional proofs on average in the full cohort and 1.86 in the capped empty-code cohort when both issuers share an epoch. The choice of cohort therefore changes the estimated magnitude of the external cost.

\subsection{Interpretation}

Common-epoch spectral negative controls used 10,000 synthetic paths, five independent seeds per condition and either four or 64 arbitrary source labels. The adjusted Rand index ranged from -0.00018 to 0.00013 across ten runs. These controls quantify spurious clustering when issuer labels carry no information. Selective clustering at this scale has one published seed per size. Accuracy of an issuer classifier on replayed Aave paths remains unmeasured. No ground-truth issuer labels exist for these addresses. The Web PKI results remain a separate sensitivity using retained effective dates in signed CRL snapshots.
